% abstract.tex -- included from paper.tex via \input{abstract}
\begin{abstract}
  EU AI Act Article~50 requires machine-readable marking of
  AI-generated content but leaves two decisive parameters -- false
  positive rate and text-quality cost -- to be chosen opaquely by
  deployers rather than specified by regulators. We present
  \emph{Markov Chain Lock} (MCL) watermarking, an active watermarking
  scheme that partitions the vocabulary into $S$ states via SHA-256
  hashing and constrains generation to follow a clockwork transition
  $s \to (s+1) \bmod S$ at selected (low-entropy) positions. Detection
  is model-free, runs in $O(n)$ hash operations, and admits
  Hoeffding-style exponential false positive bounds -- making
  third-party audit feasible without LLM access. The central result is
  a calibration formula
  $S^\star = \lceil 4 z_\alpha^2 / (\rho^2 (n-1)) + 1 \rceil$ that
  maps a target FPR $\alpha$, text length $n$, and watermark fraction
  $\rho$ to a minimum state count, turning Article~50 compliance from
  an opaque deployer choice into an auditable policy-to-configuration
  map. We additionally prove a universal robustness threshold
  $\delta^\star = 1 - 1/\sqrt{2} \approx 29.3\%$ independent of
  $(S, \rho, n)$, derive a Neyman--Pearson optimal entropy-weighted
  detector, and formalize a $1/S$ self-healing surplus under
  adversarial modification. A critical experiment tests whether
  low-entropy tokens survive paraphrasing at higher rates than
  high-entropy ones -- inverting the SWEET/EWD heuristic -- and
  decides whether the entropy gate is purely a quality improvement or
  also a robustness one.
\end{abstract}
